%% kapa: exam=1 round=23 year=2012 period=1 form=A subject=economics subject_ko=경제원론 seq=9 number=49
%% source: Q-Net 감정평가사 기출문제 (cst00309) artlSeq=5052578 file=raw/1008/5052578/2012년 제23회 감정평가사 제1차시험 기출문제(1교시 A형).pdf
%% conversion: kordoc 4.24.51 -> tools/gen1.py (auto); review=proofread-sonnet; flags=none
\begin{question}{49}
노동과 여가의 선택에 관한 설명으로 옳지 않은 것은? (단, 여가는 정상재이다.)
\begin{choices}
  \item 시간당 임금은 여가 한 시간의 기회비용이다.
  \item 시간당 임금이 상승할 경우, 소득효과만 고려하면 노동공급이 감소한다.
  \item 시간당 임금에 대해 근로소득세율이 상승할 경우, 노동공급이 항상 감소한다.
  \item 비근로소득이 증가하는 경우, 노동공급이 항상 감소한다.
  \item 시간당 임금이 상승할 경우, 대체효과만 고려하면 노동공급이 항상 증가한다.
\end{choices}
\end{question}
\answer{3}
